\documentclass[a4paper]{article}

\usepackage{ctex}
\usepackage{geometry}
\usepackage{graphicx}
\usepackage{grffile}
\usepackage{xcolor}
\usepackage{hyperref}
\usepackage{fancyhdr}
\usepackage{lastpage}
\usepackage{float}
\usepackage{booktabs}
\usepackage{longtable}
\usepackage{array}
\usepackage{tabularx}
\usepackage{enumitem}
\usepackage{caption}
\usepackage{multirow}

\geometry{a4paper,left=2.3cm,right=2.3cm,top=2.7cm,bottom=2.7cm}
\graphicspath{{./}{report_assets/}}
\hypersetup{colorlinks=true, linkcolor=black, urlcolor=blue!60!black}
\setlength{\textfloatsep}{8mm}
\setlength{\tabcolsep}{4pt}
\setlist{nosep}
\renewcommand{\arraystretch}{1.22}

\pagestyle{fancy}
\fancyhf{}
\lhead{\kaishu \leftmark}
\rhead{\kaishu 软件工程课程设计报告}
\cfoot{\thepage\ of \pageref{LastPage}}
\renewcommand{\headrulewidth}{0.1pt}
\renewcommand{\footrulewidth}{0pt}

\captionsetup{font=small, labelfont=bf}
\renewcommand{\contentsname}{目\hspace{1em}录}
\renewcommand{\figurename}{图}
\renewcommand{\tablename}{表}
\renewcommand{\today}{\number\year 年\number\month 月\number\day 日}
\renewcommand{\thefigure}{\thesection.\arabic{figure}}
\renewcommand{\thetable}{\thesection.\arabic{table}}

\newcommand{\resultfigure}[3]{
  \begin{figure}[H]
    \centering
    \IfFileExists{#1}{
      \includegraphics[width=0.92\linewidth,height=0.62\textheight,keepaspectratio]{#1}
    }{
      \fbox{\parbox{0.84\linewidth}{\centering 图片文件未找到：\texttt{\detokenize{#1}}\\[0.4em]对应内容：#2}}
    }
    \caption{#2}
    \label{#3}
  \end{figure}
}

\newcolumntype{Y}{>{\raggedright\arraybackslash}X}
\newcolumntype{P}[1]{>{\raggedright\arraybackslash}p{#1}}

\begin{document}

\begin{titlepage}
  \begin{center}
    \vspace*{1.5cm}
    {\Huge\bfseries\kaishu 软件工程课程作业}\\[0.9cm]
    {\Huge\bfseries\kaishu 软件设计报告}\\[2.2cm]
    {\LARGE\bfseries 校缘聚（CampusGather）桌游/剧本杀组局平台}\\[1.2cm]
    {\Large 团队序号：第 1 组}\\[1.6cm]

    \begin{tabular}{>{\bfseries}p{3.0cm}p{8.0cm}}
      项目名称 & 校缘聚（CampusGather）桌游/剧本杀组局平台 \\
      设计依据 & 《校缘聚需求分析文档》与本次软件设计报告要求 \\
      交付内容 & 软件设计报告、UML/ER/架构图、可运行原型、JSON 数据文件与测试 \\
    \end{tabular}

    \vspace{1.8cm}
    \begin{tabular}{ccc}
      \toprule
      姓名 & 学号 & 组内说明 \\
      \midrule
      刘砚桐 & 2313983 & 学生端与组局业务设计 \\
      李佳璞 & 2314007 & 架构、后端与报告整合 \\
      苏雨辰 & 2313828 & 需求校对与用户流程设计 \\
      史傅冠华 & 2311987 & 管理端、投诉与信用设计 \\
      朱乐晨 & 2312194 & 场地管理与数据库设计 \\
      \bottomrule
    \end{tabular}

    \vfill
    {\Large \today}
  \end{center}
\end{titlepage}

\clearpage
\tableofcontents
\newpage

\section{系统总体设计}

\subsection{项目背景与目标}

校缘聚（CampusGather）面向校园桌游、剧本杀等线下组局场景。需求分析阶段已经明确，当前学生主要依赖微信群、QQ群、贴吧等松散渠道发起活动，常见问题包括信息分散、参与者身份难以确认、活动变更通知不及时、临时跳车难追溯以及校内活动空间申请流程不透明等。桌游与剧本杀又具有固定人数、固定时长、线下协作强和信用约束强的特点，因此需要一个能够把“发布、招募、审核、场地、通知、评价、投诉”串联起来的系统。

本系统的设计目标如下：

\begin{enumerate}
  \item 为普通学生提供组局浏览、筛选、报名、退出、互评、投诉与个人信用管理能力；
  \item 为发起人提供组局发布、成员审核、活动状态维护和场地预约申请能力；
  \item 为系统管理员提供实名认证、游戏库、违规内容、投诉、信用和日志统计管理能力；
  \item 为场地/社团管理员提供场地维护、开放时段配置、预约审核、占用追踪和场地通知能力；
  \item 在课程周期内完成可运行原型，同时预留学校统一身份认证、微信小程序和通知推送等后续扩展点。
\end{enumerate}

\subsection{系统功能概述}

\begin{table}[H]
  \centering
  \begin{tabularx}{\linewidth}{P{2.5cm} P{5.4cm} Y}
    \toprule
    用户域 & 核心功能 & 设计说明 \\
    \midrule
    普通学生 & 浏览/筛选组局、查看详情、申请加入、退出、互评、投诉、个人资料（含头像与标签）、信用记录、实名认证提交 & 未认证用户只能浏览公开信息；认证后才能参与影响活动与信用的操作；可通过个人主页编辑头像、昵称、标签等。 \\
    发起人 & 发布组局、编辑组局、审核报名、取消组局、发起场地申请 & 发起人是普通学生在特定组局中的业务角色，权限限定在自己创建的组局。 \\
    系统管理员 & 实名认证审核、账号状态管理、用户列表查看、游戏库、违规内容、投诉与信用、通知、日志统计 & 后台接口统一走角色鉴权，所有关键操作写入管理员操作日志；可查看全部用户列表并对用户进行认证审核与账号状态变更。 \\
    场地管理员 & 场地信息、开放时段、预约审核、占用记录、场地通知 & 只能处理被授权管理的场地，审核结果同步影响组局场地状态。 \\
    外部系统 & 学校统一身份认证、通知推送、日志监控 & 当前原型采用模拟适配器；部署阶段替换为真实接口。 \\
    \bottomrule
  \end{tabularx}
  \caption{系统功能概览}
\end{table}

\subsection{系统总体架构设计}

系统采用客户端--服务器与分层 MVC 相结合的架构。客户端层包含微信小程序学生端和管理后台；应用服务层使用 Node.js REST API 承载鉴权、组局、场地、信用、投诉、通知等模块；数据层统一采用 JSON 文件持久化，便于本地演示、版本追踪和课程验收。

\resultfigure{architecture.png}{系统总体架构设计}{fig:architecture}

\begin{table}[H]
  \centering
  \begin{tabularx}{\linewidth}{P{2.4cm} P{6.1cm} Y}
    \toprule
    层次 & 职责 & 可实践约束 \\
    \midrule
    表现层 & 移动端与后台页面，负责表单录入、状态展示、操作确认、错误提示 & 前端不直接访问数据文件；所有业务操作经 REST API 完成。 \\
    接口/控制层 & 路由分发、参数校验、登录态解析、角色权限校验、统一响应格式 & 错误码统一为 VALIDATION\_ERROR、FORBIDDEN、NOT\_FOUND、CONFLICT 等。 \\
    业务服务层 & 组局、报名、场地、投诉、信用、通知、日志等业务规则 & 关键状态变更集中在服务层，避免前端绕过规则。 \\
    数据访问层 & 通过 JSONStore 封装 CRUD、查询、保存与种子数据 & 每个实体集合独立保存为 JSON 文件，写入时使用临时文件替换。 \\
    基础设施层 & 学校认证、推送、监控、备份 & 均使用适配器模式，原型可模拟，部署可替换。 \\
    \bottomrule
  \end{tabularx}
  \caption{架构分层说明}
\end{table}

\subsection{技术选型说明}

\begin{longtable}{P{2.3cm} P{4.1cm} P{8.8cm}}
  \toprule
  类别 & 选型 & 理由 \\
  \midrule
  \endfirsthead
  \toprule
  类别 & 选型 & 理由 \\
  \midrule
  \endhead
  前端 & 微信小程序/响应式 Web 原型 & 符合校园移动使用场景；Web 原型便于课程验收直接运行。 \\
  后端 & Node.js 原生 HTTP + 模块化服务 & 异步 I/O 适合高并发浏览与报名；无外部依赖，便于离线演示和部署。 \\
  接口 & RESTful API + JSON & 易于小程序调用、调试和生成接口文档；统一错误码便于前后端协作。 \\
  数据存储 & JSON 文件持久化 & 不依赖外部数据库服务，降低演示环境成本；集合结构由 seed.js 和运行数据共同体现。 \\
  图表 & Graphviz DOT & 本地可生成 UML、ER 与架构图，源文件可追踪，后续可转 PlantUML 或 draw.io。 \\
  测试 & Node.js 内置 node:test & 不依赖网络安装包即可覆盖核心接口与业务规则。 \\
  部署 & Node.js + 静态前端 + JSON 数据目录 & 结构简单，适合课程验收与本地演示；后续可通过缓存、备份和访问控制增强稳定性。 \\
  \bottomrule
  \caption{技术选型说明}
\end{longtable}

\section{详细设计}

\subsection{模块划分与模块功能设计}

\begin{longtable}{P{3.0cm} P{4.2cm} P{4.0cm} P{4.0cm}}
  \toprule
  模块 & 主要职责 & 输入/输出 & 关键规则 \\
  \midrule
  \endfirsthead
  \toprule
  模块 & 主要职责 & 输入/输出 & 关键规则 \\
  \midrule
  \endhead
  Auth/User & 登录、当前用户、个人资料（含头像编辑）、实名认证提交与管理员审核、用户列表管理、角色与账号状态变更 & 输入学号或用户编号；输出用户视图（含头像URL、认证提交信息、权限） & 未认证用户不可发布、报名、评价、投诉；学号和联系方式仅本人和管理员可见；头像限定预设选项。 \\
  GameLib & 维护桌游/剧本库，供发布和筛选使用 & 游戏名称、类型、人数、时长、难度、标签 & 下架条目不可被新组局选择。 \\
  Session & 组局发布、列表筛选、详情、编辑、取消、完结 & 活动时间、地点、人数、信用要求、加入方式 & 时间必须合法；名额不得超上限；关键变更发通知。 \\
  Application/Member & 报名申请、发起人审核、退出、成员名单 & 申请备注、审核动作、退出原因 & 审核制先进入待审核；直接加入走容量和冲突校验。 \\
  Venue & 场地资源、开放时段、预约申请、审核、占用记录 & 场地、时间、申请说明、审核意见 & 容量、开放状态、时间冲突和管理权限必须校验。 \\
  Review/Credit & 活动互评、信用分流水、跳车或投诉扣分 & 评分、评价内容、信用变更原因 & 仅实际成员可评价；信用变更必须有业务关联。 \\
  Complaint & 投诉提交、受理、驳回/成立、处理结果 & 投诉原因、证据、目标用户 & 管理员处理后同步通知双方并写日志。 \\
  Notification/Log/Stats & 站内通知、已读状态、管理员操作日志、平台统计 & 业务事件和操作上下文 & 日志不可由普通接口删除；统计仅管理员可见。 \\
  \bottomrule
  \caption{模块功能设计}
\end{longtable}

\subsection{模块调用关系}

控制器只负责请求解析和响应封装，不直接改写业务数据。每个控制器调用对应服务；服务之间通过显式方法协作。例如报名成功后，SessionService 调用 NotificationService 生成通知，并调用 LogService 写入操作记录；投诉成立后，ComplaintService 调用 CreditService 写入信用流水，再通知投诉双方。

\begin{table}[H]
  \centering
  \begin{tabularx}{\linewidth}{P{2.7cm} P{3.2cm} P{5.2cm} Y}
    \toprule
    入口操作 & 主服务 & 协作服务 & 数据实体 \\
    \midrule
    发布组局 & SessionService & GameLibService、VenueService、NotificationService、LogService & game\_sessions、session\_members、venue\_reservations、notifications、admin\_logs \\
    申请加入 & SessionService & UserService、CreditService、NotificationService & session\_applications、session\_members、credit\_records \\
    场地审核 & VenueService & SessionService、NotificationService、LogService & venue\_reservations、venues、game\_sessions \\
    处理投诉 & ComplaintService & CreditService、NotificationService、LogService & complaints、credit\_records、notifications、admin\_logs \\
    实名认证审核 & UserService & ExternalAuthAdapter、NotificationService、LogService & users、notifications、admin\_logs \\
    实名认证提交 & UserService & — & users \\
    \bottomrule
  \end{tabularx}
  \caption{模块调用关系}
\end{table}

\subsection{核心业务流程设计}

发布组局流程：用户认证通过后进入发布页面，选择游戏库条目并填写标题、人数、时间、场地、加入方式和最低信用分。系统校验时间合法性、人数范围、游戏状态和场地容量与时间冲突；场地选择为必填项，发布组局时自动创建已通过的场地预约记录（venue\_status 置为 approved），组局地点字段自动由场地信息生成。保存组局后将发起人写入成员表。后续发起人仍可通过编辑组局更换场地或调整时间，系统自动同步场地预约。

实名认证流程：学生进入个人主页，填写真实姓名、学号和联系方式提交认证申请，auth\_status 变为 pending。管理员在后台用户列表中查看待审核用户及其提交信息，选择通过（approve）或驳回（reject）并填写审核意见。通过后用户 auth\_status 变为 verified，获得发布、报名、评价等操作权限；驳回后用户可修改信息重新提交。审核结果通过通知模块反馈给申请人。

报名与成员审核流程：学生提交申请后，系统检查认证状态、账号状态、信用分、组局状态、名额、重复报名和时间冲突。直接加入模式立即生成成员；审核制生成待审核申请，由发起人通过或拒绝。所有结果通过通知模块反馈。

\resultfigure{sequence_join.png}{申请加入组局顺序图}{fig:sequence-join}

场地审核流程：发起人提交预约后，场地服务校验场地状态、容量和时间范围，并检查已通过预约是否冲突。场地管理员只能审核自己管理的场地。通过后预约状态变为 approved，组局详情显示场地已确认；驳回时保存原因并通知发起人。

\resultfigure{sequence_venue.png}{场地预约审核顺序图}{fig:sequence-venue}

\subsection{接口设计}

\begin{longtable}{P{2.6cm} P{5.4cm} P{2.2cm} P{5.8cm}}
  \small
  \toprule
  方法 & 路径 & 角色 & 说明 \\
  \midrule
  \endfirsthead
  \toprule
  方法 & 路径 & 角色 & 说明 \\
  \midrule
  \endhead
  GET & \texttt{/api/health} & 公开 & 健康检查与版本信息 \\
  POST & \texttt{/api/auth/login} & 公开 & 使用学号/角色演示登录，返回用户视图 \\
  GET & \texttt{/api/users/me} & 登录用户 & 查看个人资料与权限 \\
  PUT & \texttt{/api/users/me} & 登录用户 & 编辑个人资料（昵称、头像、标签、联系方式） \\
  POST & \texttt{/api/users/me/auth} & 学生 & 提交实名认证申请（真实姓名、学号、联系方式） \\
  GET & \texttt{/api/users/me/credit} & 登录用户 & 查看个人信用分与信用流水 \\
  GET & \texttt{/api/users} & 管理员 & 查看全部用户列表 \\
  PATCH & \texttt{/api/users/:id/auth} & 管理员 & 审核用户实名认证（通过/驳回） \\
  PATCH & \texttt{/api/users/:id/status} & 管理员 & 变更用户账号状态（active/muted/limited/banned） \\
  GET\allowbreak/\allowbreak POST & \texttt{/api/games} & 公开\allowbreak/\allowbreak 管理员 & 查询游戏库；管理员新增游戏 \\
  PATCH & \texttt{/api/games/:id} & 管理员 & 更新游戏条目信息 \\
  GET\allowbreak/\allowbreak POST & \texttt{/api/sessions} & 公开\allowbreak/\allowbreak 认证学生 & 筛选组局；发布组局（必选场地） \\
  GET & \texttt{/api/sessions/mine} & 登录用户 & 查看我发布或参与的组局 \\
  GET\allowbreak/\allowbreak PATCH & \texttt{/api/sessions/:id} & 公开\allowbreak/\allowbreak 发起人 & 查看详情（含场地与成员信息）；发起人编辑组局 \\
  POST & \texttt{/api/sessions/\allowbreak :id/\allowbreak applications} & 认证学生 & 申请加入组局或直接加入 \\
  PATCH & \texttt{/api/applications/:id} & 发起人 & 通过/拒绝报名申请 \\
  POST & \texttt{/api/sessions/:id/leave} & 成员 & 退出组局并按规则记录信用影响 \\
  POST & \texttt{/api/sessions/:id/cancel} & 发起人 & 取消组局 \\
  POST & \texttt{/api/sessions/:id/finish} & 发起人 & 结束组局 \\
  POST & \texttt{/api/sessions/:id/reviews} & 成员 & 活动结束后互评 \\
  GET\allowbreak/\allowbreak POST\allowbreak/\allowbreak PATCH & \texttt{/api/complaints} & 学生\allowbreak/\allowbreak 管理员 & 提交、查询、处理投诉 \\
  GET\allowbreak/\allowbreak POST & \texttt{/api/venues} & 场地管理员 & 查询场地列表；新增场地 \\
  GET\allowbreak/\allowbreak PATCH\allowbreak/\allowbreak DELETE & \texttt{/api/venues/:id} & 场地管理员 & 查看、编辑、删除场地（删除时级联取消关联组局） \\
  GET\allowbreak/\allowbreak POST\allowbreak/\allowbreak PATCH & \texttt{/api/venue-reservations} & 发起人\allowbreak/\allowbreak 场地管理员 & 提交、查询、审核场地预约 \\
  GET\allowbreak/\allowbreak PATCH & \texttt{/api/notifications} & 登录用户 & 查看通知与标记已读 \\
  GET & \texttt{/api/admin/logs} & 管理员 & 查看操作日志 \\
  GET & \texttt{/api/admin/stats} & 管理员 & 查看平台统计 \\
  \bottomrule
  \caption{主要 REST 接口设计}
\end{longtable}

统一响应格式为：成功时返回 \texttt{\{ success: true, data \}}；失败时返回 \texttt{\{ success: false, error: \{ code, message, details \} \}}。客户端根据 HTTP 状态码和 code 显示明确错误，例如名额已满、权限不足、信用分不足、场地冲突、重复报名等。

\section{UML 设计}

\subsection{用例图}

\resultfigure{use_case.png}{系统用例图}{fig:use-case}

用例图覆盖普通学生、发起人、系统管理员、场地/社团管理员与学校统一身份认证接口。发起人不是独立账号类型，而是普通学生在自己发布的组局中的临时角色。

\subsection{类图}

\resultfigure{class_diagram.png}{核心领域类图}{fig:class}

类图以业务领域对象为中心。User 与 GameSession、SessionApplication、SessionMember 共同支撑组局；Venue 与 VenueReservation 支撑场地审批；Review、Complaint、CreditRecord 共同支撑信用闭环；Notification 与 AdminLog 为横切能力。

\subsection{顺序图}

核心顺序图已在第 2 章给出，分别展示申请加入组局与场地预约审核。实现时应保证顺序图中的校验步骤在服务层执行，而不是只在前端提示。

\subsection{部署图}

\resultfigure{deployment.png}{系统部署图}{fig:deployment}

部署设计以课程原型为核心形态：一个 Node.js 进程同时提供静态页面、REST API 和 JSON 文件持久化能力；正式试点时可在外层接入 HTTPS、日志告警、学校统一身份认证和推送服务。

\section{数据存储设计}

\subsection{数据实体关系图}

\resultfigure{er_diagram.png}{数据实体关系图}{fig:er}

数据模型以 users 为主体，game\_sessions 为活动核心。报名申请、成员、预约、评价、投诉、信用、通知、日志均通过 id 字段关联到用户或组局，便于追溯完整业务链路。

\subsection{主要数据集合设计}

\begin{longtable}{P{3.2cm} P{5.4cm} P{4.0cm} P{3.0cm}}
  \toprule
  集合名 & 关键字段 & 主键/关联 & 说明 \\
  \midrule
  \endfirsthead
  \toprule
  集合名 & 关键字段 & 主键/关联 & 说明 \\
  \midrule
  \endhead
  users & id, student\_no, name, nickname, role, auth\_status, credit\_score, status, avatar, tags, contact, auth\_submission\_name, auth\_submission\_student\_no, auth\_submission\_contact, auth\_submission\_note, auth\_submitted\_at, auth\_review\_reason, auth\_reviewed\_at, status\_reason & PK id；UNIQUE student\_no & 用户、管理员与场地管理员统一建模，通过 role 区分权限。ID 采用 5 位数字编码：10001--11000 为管理员，11001--99999 为普通用户。avatar 存储头像文件名，auth\_submission 系列字段支持实名认证提交流程。 \\
  game\_libs & id, name, type, min\_players, max\_players, duration\_minutes, difficulty, description, tags, status & PK id & 桌游/剧本库，供发布组局和筛选使用。 \\
  game\_sessions & id, host\_id, game\_id, title, description, start\_time, end\_time, location, max\_members, current\_members, min\_credit\_required, join\_mode, status, venue\_status, cancel\_reason & PK id；FK host\_id；FK game\_id & 一次具体组局活动。venue\_status 记录场地预约状态，cancel\_reason 记录取消原因。 \\
  session\_applications & id, session\_id, applicant\_id, message, status, apply\_time, review\_time, review\_reason & PK id；FK session\_id；FK applicant\_id & 审核制报名申请与直接加入记录。 \\
  session\_members & id, session\_id, user\_id, member\_role, join\_time & PK id；FK session\_id；FK user\_id；UNIQUE(session\_id,user\_id) & 已加入成员名单。 \\
  venues & id, name, location, capacity, manager\_id, available\_time, open\_rules, status, description & PK id；FK manager\_id & 校园场地基础数据与开放规则。 \\
  venue\_reservations & id, venue\_id, session\_id, applicant\_id, reviewer\_id, start\_time, end\_time, status, review\_reason & PK id；FK venue\_id；FK session\_id；FK applicant\_id & 场地预约审核记录。 \\
  reviews & id, session\_id, reviewer\_id, target\_user\_id, score, content, created\_at & PK id；FK session\_id；FK reviewer/target & 活动结束后互评。 \\
  complaints & id, reporter\_id, target\_user\_id, session\_id, reason, evidence, status, result, handled\_by & PK id；FK reporter/target；FK session\_id & 投诉受理、处理和结果记录。 \\
  credit\_records & id, user\_id, session\_id, complaint\_id, change\_value, reason, created\_at & PK id；FK user\_id；FK session\_id；FK complaint\_id & 信用分变化流水。 \\
  notifications & id, user\_id, type, title, content, related\_type, related\_id, read\_at, created\_at & PK id；FK user\_id & 站内通知与后续微信订阅消息记录。 \\
  admin\_logs & id, operator\_id, action, object\_type, object\_id, result, remark, created\_at & PK id；FK operator\_id & 后台关键操作审计日志。 \\
  \bottomrule
  \caption{主要数据表设计}
\end{longtable}

\subsection{关键查询与一致性设计}

\begin{enumerate}
  \item \texttt{game\_sessions} 按 \texttt{status}、\texttt{start\_time}、\texttt{game\_id} 和 \texttt{host\_id} 进行内存筛选与排序，支撑列表和个人组局查询。
  \item \texttt{session\_members} 通过服务层检查 \texttt{(session\_id, user\_id)} 组合，防止重复报名；报名申请按 \texttt{status} 区分待审核记录。
  \item \texttt{venue\_reservations} 在服务层按 \texttt{venue\_id}、时间区间和状态进行冲突检测，避免同一场地被重复占用。
  \item 报名审批、退出扣分、投诉成立后的信用变更由单进程同步写入 JSON 文件，保证课程演示场景下的数据一致性。
\end{enumerate}

\section{安全、性能与可扩展性设计}

\subsection{安全性设计}

\begin{table}[H]
  \centering
  \begin{tabularx}{\linewidth}{P{3.0cm} Y}
    \toprule
    安全点 & 设计方案 \\
    \midrule
    身份认证 & 所有写操作必须登录；发布、报名、评价、投诉要求 \texttt{auth\_status=verified}。学校统一身份认证采用适配器接入，原型可模拟审核。 \\
    权限控制 & 普通学生、系统管理员、场地管理员按角色授权；发起人只能管理自己发布的组局；场地管理员只能审核自己管理的场地。 \\
    隐私保护 & 学号、联系方式、投诉证据等敏感字段默认不在公开接口返回；对外展示信用记录采用次数和类型脱敏。 \\
    输入校验 & 标题、简介、评价、投诉内容进行长度、必填和敏感词校验；人数、时间、评分等使用类型和范围校验。 \\
    审计日志 & 实名认证、账号状态、违规处理、投诉处理、场地审核、信用调整等后台动作写入 AdminLog，至少保留 6 个月。 \\
    \bottomrule
  \end{tabularx}
  \caption{安全性设计}
\end{table}

\subsection{性能优化设计}

\begin{enumerate}
  \item 列表接口支持分页和条件筛选，首页只返回必要摘要，详情页再加载完整信息。
  \item 报名和审核接口采用原子校验：先查状态、名额、重复成员和时间冲突，再写入成员/申请。
  \item 热门查询字段建立索引，游戏库和场地基础数据可在客户端短期缓存。
  \item 通知发送采用先落库、后推送策略；推送失败不影响主业务完成，可重试。
  \item 目标指标沿用需求文档：列表与筛选 2 秒内返回，关键操作 3 秒内反馈，日常支持 100 名用户同时浏览。
\end{enumerate}

\subsection{可扩展性设计}

\begin{enumerate}
  \item 活动类型以 \texttt{game\_libs.type} 和 tags 扩展，后续可加入狼人杀、密室、社团活动等。
  \item 角色权限集中在鉴权中间件与用户 role 字段，后续可加入 DM、社团负责人、商家等角色。
  \item 推荐算法独立为 RecommendationService，初期基于时间/类型/地点筛选，后续可引入兴趣标签和历史参与记录。
  \item 外部学校认证、微信订阅消息、短信或邮件都通过 Adapter 封装，避免业务服务直接依赖第三方 SDK。
  \item JSONStore 集中封装读写逻辑，业务服务只依赖集合接口，便于后续扩展备份、校验和导入导出能力。
\end{enumerate}

\subsection{异常处理机制}

\begin{table}[H]
  \centering
  \begin{tabularx}{\linewidth}{P{3.0cm} P{5.9cm} Y}
    \toprule
    异常类型 & 处理策略 & 用户提示 \\
    \midrule
    参数不完整/格式错误 & 控制层返回 400 与字段级 details & 指出具体字段，例如“活动结束时间必须晚于开始时间”。 \\
    权限不足 & 鉴权层返回 403 并记录异常访问 & 提示“当前账号无权执行该操作”。 \\
    资源不存在 & 服务层返回 404 & 提示目标组局、场地或申请不存在。 \\
    业务冲突 & 服务层返回 409 & 提示名额已满、重复报名、时间冲突或场地冲突。 \\
    外部接口失败 & 适配器超时重试并降级为待人工审核 & 提示认证/推送暂不可用，稍后重试或等待管理员处理。 \\
    \bottomrule
  \end{tabularx}
  \caption{异常处理机制}
\end{table}

\section{实现对应关系与测试计划}

\subsection{报告设计与代码实现对应关系}

\begin{table}[H]
  \centering
  \begin{tabularx}{\linewidth}{P{3.0cm} P{5.4cm} Y}
    \toprule
    报告模块 & 代码实现 & 验收方式 \\
    \midrule
    接口/控制层 & \texttt{src/server.js}, \texttt{src/router.js} & 访问 \texttt{/api/health}，使用前端或测试脚本调用核心接口。含头像静态资源路由 \texttt{/profile\_photo/}。 \\
    业务服务层 & \texttt{src/services/userService.js}, \texttt{sessionService.js}, \texttt{venueService.js}, \texttt{complaintService.js} 等 & 通过 \texttt{node --test} 覆盖登录、认证提交与审核、发布（含场地绑定）、报名、审核、场地冲突、投诉信用变更。 \\
    数据访问层 & \texttt{src/database/jsonStore.js}, \texttt{src/database/seed.js} 与 \texttt{data/*.json} & 重启后数据不丢失；seed.js 预置 6 个用户（含头像、认证状态）、2 个场地、3 款游戏和 2 条组局示例。 \\
    数据模型设计 & \texttt{src/database/seed.js}, \texttt{data/*.json} & 检查集合字段与本报告一致；users 集合含 avatar 与 auth\_submission 系列字段。 \\
    前端演示 & \texttt{public/index.html}, \texttt{public/app.js}, \texttt{public/styles.css} & 浏览器打开本地服务即可完成学生个人主页编辑、实名认证提交、管理员审核、场地选择发布组局、参与者详情查看等完整流程。 \\
    静态资源 & \texttt{data/profile\_photo/}（12 张头像图片） & 头像通过 \texttt{/profile\_photo/文件名} 访问，支持 default、adm、1--10 共 12 个选项。 \\
    \bottomrule
  \end{tabularx}
  \caption{设计与实现映射}
\end{table}

\subsection{原型边界}

本次课程实现以完整业务闭环为目标，不接入真实微信 AppID、学校统一身份认证和线上推送密钥；这些环境参数需要学校或课程平台确认后替换。原型统一使用 JSON 文件持久化，便于课堂演示、数据重置和版本检查。

\subsection{测试计划}

\begin{table}[H]
  \centering
  \begin{tabularx}{\linewidth}{P{3.0cm} P{6.0cm} Y}
    \toprule
    阶段 & 内容 & 完成标准 \\
    \midrule
    单元/接口测试 & 健康检查、登录、筛选、发布（含场地绑定）、申请、审核、退出、取消/完结、投诉、实名认证提交与审核、场地审核与删除 & \texttt{npm test} 全部通过。 \\
    业务验收 & 学生端完成“编辑个人主页--实名认证提交--浏览--报名--评价/投诉”；管理员端完成“认证审核/用户列表/投诉/内容库”；场地端完成“场地管理/预约审核”。 & 前端页面可演示，头像上传与展示、ID 分配方案、数据持久化可追溯。 \\
    性能检查 & 构造列表筛选和连续报名请求 & 关键接口在本地 3 秒内响应，无超额报名。 \\
    安全检查 & 未登录、未认证、越权账号分别调用写接口 & 返回 401/403，不修改数据。 \\
    提交 & 报告源码、图片、PDF、源码、README、JSON 数据说明、测试 & 推送到 GitHub main 分支。 \\
    \bottomrule
  \end{tabularx}
  \caption{测试与交付计划}
\end{table}

\end{document}
